%----------------------------------------------------------------------------------------
% БҮЛЭГ 1: ОРШИЛ
%----------------------------------------------------------------------------------------

\chapter{Оршил}

\section{Судалгааны үндэслэл}

Орчин үед дэлхий даяар видео тандалтын систем хурдацтай өргөжиж байна. IHS Markit-ийн 2019 оны тайланд дэлхийд ашиглагдаж байгаа хяналтын камерын тоо 770 саяд хүрсэн бөгөөд 2025 он гэхэд 1 тэрбумаас давах төлөвтэй байгааг тодорхойлсон \cite{IHSMarkit2019}. Comparitech-ийн судалгаагаар Хятад улс нэг хүнд ноогдох камерын тоогоор тэргүүлж, 100 хүнд 54.64 камер ноогддог бол Монгол улсад ч энэ чиг хандлага нэмэгдэж байна \cite{Comparitech2020}.

Хүн амын төвлөрөл, олон нийтийн орчны идэвхжил нэмэгдэхийн хэрээр аюулгүй байдлыг хангах асуудал улам бүр чухал болж байна. Энэ хүрээнд видео тандалтын систем нь гэмт хэрэг, зөрчлийг илрүүлэх, урьдчилан сэргийлэх, нотлох баримт бүрдүүлэх үндсэн технологийн шийдэл болон өргөн хэрэглэгдэж байна. Олон улсын судалгаануудаас харахад хяналтын камерын систем нь гэмт хэргийн илрүүлэлт, хяналтын үр нөлөөг нэмэгдүүлэхэд чухал үүрэгтэй болох нь тогтоогдсон \cite{piza2019cctv}.

Хяналтын камеруудыг олон нийтийн газар, дотуур байр, боловсролын байгууллага, худалдааны төв, тээврийн хөдөлгөөн их төвлөрдөг бүс зэрэг төрөл бүрийн орчинд байршуулснаар тасралтгүй видео мэдээлэл цуглуулж, тухайн орчны нөхцөл байдлыг хянах боломжийг бүрдүүлдэг. Эдгээр бичлэг нь зөрчил, гэмт хэргийн дараах дүн шинжилгээ хийх, үйл явдлын нөхцөл байдлыг сэргээн тогтоох, холбогдох байгууллагын шийдвэр гаргалтад дэмжлэг үзүүлэх чухал мэдээллийн эх үүсвэр болдог.

Гэвч одоогийн ихэнх видео тандалтын систем нь хүний оператороор хянагддаг уламжлалт хэлбэртэй хэвээр байна. Нэг оператор нэгэн зэрэг 4--16 дэлгэц хянах боломжтой боловч камерын тоо нэмэгдэхийн хэрээр бүх бичлэгийг тасралтгүй хянах нь хүний нөөцийн хувьд хязгаарлагдмал бөгөөд анхаарал сулрах, алдаа гаргах, зөрчлийг цаг тухайд нь илрүүлэхгүй байх эрсдэлийг нэмэгдүүлдэг. Урт хугацааны хяналтын явцад хүний анхаарал мэдэгдэхүйц буурдаг болохыг олон судалгаа нотолсон \cite{mackworth1948vigilance,warm2008vigilance}.

Сүүлийн жилүүдэд хиймэл оюун ухаан, ялангуяа компьютерийн хараа болон VLM-ийн хөгжил нь видео тандалтын системийг автоматжуулах шинэ боломжийг нээж байна. Эдгээр загварууд нь зөвхөн дүрсийг таних төдийгүй тухайн үйл явдлын агуулгыг ойлгож, текстэн тайлбар үүсгэх чадвартайгаараа уламжлалт компьютерийн хараа аргуудаас давуу талтай юм.

Монгол улсын тухайд дотоодын аюулгүй байдлын байгууллагууд хяналтын камерын бичлэгийг Монгол хэл дээр тайлбарлах, дүн шинжилгээ хийх хэрэгцээтэй тулгардаг. Гэвч одоогоор зодооны зөрчлийг Монгол хэлээр тайлбарлах, автоматаар илрүүлэх систем хөгжөөгүй байна. Энэ нь нэн тэргүүний шаардлага болж байгаа бөгөөд энэхүү судалгаа нь тухайн хоосон орон зайг нөхөхөд чиглэнэ.

\section{Асуудлын тодорхойлолт}

Хяналтын камерын тоо нэмэгдэхийн хэрээр бодит цагийн хяналт нь хүний нөөцийн хувьд боломжгүй болж байна. Үүнийг шийдвэрлэх автоматчилсан систем шаардлагатай боловч одоогийн шийдлүүдэд дараах гол асуудлууд байна:

\begin{itemize}
    \item \textbf{Тайлбарын дутагдал:} Уламжлалт компьютерийн хараа загварууд зодооны үйл явдлыг ``хүчирхийлэлтэй'' эсвэл ``хүчирхийлэлгүй'' гэж ангилдаг боловч яагаад, хэрхэн болсон талаар тайлбар өгдөггүй. Хэрэг шалгах байгууллагуудад хуулийн нотлох баримт болох нарийвчилсан тайлбар шаардлагатай.
        \item \textbf{Монгол хэлний дутагдал:} Одоо байгаа бүх систем Монгол хэлний тайлбар гаргах чадваргүй бөгөөд энэ нь дотоодын хэрэглэгчдэд ашиглахад томоохон саад болдог.
    \item \textbf{Тооцооллын хүчин чадлын хязгаарлалт:} Өндөр гүйцэтгэлтэй том загвар нь бодит цагийн орчинд ашиглахад тохиромжгүй тул хурдан, хөнгөн загвар боловсруулах шаардлагатай.
    \item \textbf{Өгөгдлийн хүрэлцээ муу байдал:} Монгол хяналтын камерын нөхцөлд тохирсон, Монгол хэлний тайлбартай сургалтын өгөгдөл байхгүй.
\end{itemize}

\section{Судалгааны зорилго}

Энэхүү судалгааны зорилго нь хяналтын камерын видео бичлэгт дүн шинжилгээ хийж, зодоон болон зөрчил зэрэг хүчирхийллийн шинжтэй үйлдлүүдийг хиймэл оюунаар илрүүлэх, илэрсэн үйл явдлын талаар Монгол хэл дээр тайлбар бүхий мэдээлэл үүсгэх, улмаар аюулгүй байдлын тасралтгүй хяналтыг илүү үр дүнтэй болгоход оршино.

\section{Судалгааны зорилт}

Дипломын ажлын хүрээнд дараах зорилтуудыг хэрэгжүүлэхээр төлөвлөж байна:

\begin{itemize}
    \item Хяналтын камерын бичлэгээс зодоон, зөрчлийн шинжтэй үйл явдлыг автоматаар илрүүлэхэд шаардлагатай өгөгдөл, аргачлалыг судлах.
    \item Олон төрлийн өгөгдөл боловсруулах том хэлний загвар ашиглан видео агуулгыг ойлгох, зөрчлийн нөхцөл байдлыг тодорхойлох болон Монгол хэлээр тайлбар үүсгэх боломжийг боловсруулах.
    \item Зодооны зөрчлийн ангилал, оролцогчдын үйлдэл, үйл явдлын дарааллыг харгалзан ангилах болон тайлбарлах загварын гүйцэтгэлийг сайжруулах.
    \item Боловсруулсан аргачлалыг туршилтын орчинд үнэлж, бодит хяналтын системд ашиглах боломжийг тодорхойлох.
\end{itemize}

\section{Судалгааны таамаглал}

Судалгааны үндсэн таамаглал нь дараах байдлаар тодорхойлогдоно. VLM загварын давуу тал нь дүрс болон текст мэдээллийг хамтад нь ойлгож, өндөр түвшний үзэгдлийн утга агуулгыг ойлгон тайлбарлах чадварт оршино. Иймд том хэмжээний VLM загварыг багш загвар болгон ашиглаж, түүнээс үүсгэсэн өндөр чанарын шошгололт нь жижиг загварын гүйцэтгэлийг мэдэгдэхүйц сайжруулах боломжтой гэж үзэж байна.

Тодруулбал: ``Qwen3-VL-30B-A3B багш загварын мэдлэгийг Qwen2.5-7B загварт LoRA нарийн тохируулгаар дамжуулах нь хяналтын камерын бичлэгээс зодооны зөрчлийг илрүүлэх болон Монгол хэлээр тайлбарлах чадварыг сайжруулах боломжтой.''

Таамаглалыг дэмжих онолын үндэслэл:

\begin{itemize}
    \item Мэдлэг дамжуулалт нь багш загварын мэдлэгийг сурагч загварт дамжуулдаг болохыг Hinton et al. \cite{hinton2015distilling} нотолсон.
    \item LoRA нь параметрийн 1--2\%-г сургаснаар бүтэн нарийн тохируулгатай дүйцэхүйц үр дүн гаргадаг болохыг Hu et al. \cite{hu2022lora} харуулсан.
    \item VLM загварууд нь видео агуулгыг ойлгоод тайлбар гаргах чадвартайг Liu et al. \cite{liu2023llava} нотолсон.
\end{itemize}

Дээрх онолын үндэслэлд тулгуурлан энэхүү судалгаанд дараах ажлуудыг хэрэгжүүлнэ:

\begin{itemize}
    \item Зодооны зөрчлийн зэрэглэл (байхгүй, бага, дунд, өндөр)-ийг тодорхойлох ангиллын систем боловсруулах.
    \item Qwen2.5-7B загварыг LoRA нарийвчилсан тохируулгын аргаар сайжруулж, 5 үеийн сургалтаар гүйцэтгэлийг нэмэгдүүлэх.
    \item Нарийн тохируулсан загварыг Qwen2-VL, Qwen2.5-VL загваруудтай харьцуулан үнэлэх.
    \item WandB тасралтгүй хяналт болон хадгалалтын цэгийн систем ашиглан сургалтын явцыг хянах, үнэлэх.
    \item Монгол хэлээр нарийвчилсан, бүтэцтэй тайлбар гаргах чадварыг туршилтаар нотлох.
\end{itemize}

\section{Хамрах хүрээ ба хязгаарлалт}

\subsection*{Хамрах хүрээ}

Энэхүү судалгаа дараах хүрээнд явагдана:

\begin{itemize}
    \item \textbf{Өгөгдөл:} RWF-2000 өгөгдлийн сан, Real Life Violence Situations өгөгдлийн багц (зодооны хэсэг), хүчирхийлэл өгөгдлийн багц.
    \item \textbf{Загвар:} Qwen загварууд (Qwen2-VL, Qwen2.5-VL, Qwen3-VL-30B-A3B).
    \item \textbf{Арга зүй:} Багш-сурагч загварын мэдлэг дамжуулалт арга, LoRA нарийн тохируулга.
    \item \textbf{Гаралт:} Зөрчил, зодооны илрүүлэлт : Монгол хэлний тайлбар.
    \item \textbf{Платформ:} NVIDIA A100 80GB GPU.
\end{itemize}

\subsection*{Хязгаарлалт}

Судалгааны хязгаарлалтуудыг дараах байдлаар тодорхойлно:

\begin{itemize}
    \item \textbf{Өгөгдлийн багцын хязгаарлалт:} Монгол хяналтын камерын бодит зодооны бичлэг хомс тул олон улсын нийтийн өгөгдлийн багц ашигласан. Монголын бодит нөхцөлд туршилт хийх боломж хязгаарлагдмал.
    \item \textbf{Тооцооллын нөөц:} Багш загвар нь өндөр тооцооллын нөөц шаарддаг тул шууд видео урсгал боловсруулах хэрэглээнд нэмэлт оновчлол хэрэгтэй.
    \item \textbf{Хэлний хязгаарлалт:} Монгол хэлний тайлбарын чанар нь суурь загварын Монгол хэлний сургалтын өгөгдлийн сангийн хэмжээтэй шууд холбоотой.
    \item \textbf{Нэгэн зэрэг олон камер:} Одоогийн систем нэгэн зэрэг олон камерын урсгалыг зэрэг боловсруулах чадваргүй.
    \item \textbf{Видео формат:} Судалгаа нь кадрт суурилсан оролтод суурилсан тул шууд видео урсгал боловсруулах хэсэгт нэмэлт хөгжүүлэлт шаардагдана.
\end{itemize}

\section{Судалгааны ач холбогдол}

\subsection*{Онолын ач холбогдол}

Энэхүү судалгааны онолын ач холбогдол нь видео тандалтын системд VLM-ийг ашиглах шинэ хандлагыг санал болгож буйд оршино. Уламжлалт видео шинжилгээний аргууд нь ихэвчлэн зөвхөн дүрсийн мэдээлэлд тулгуурлан үйл явдлыг ангилах түвшинд ажилладаг бол энэхүү судалгаа нь дүрс болон текст мэдээллийг хослуулан илүү өндөр түвшний утга агуулгын ойлголт, тайлбар үүсгэх боломжийг харуулж байна.

Ялангуяа багш-сурагч мэдлэг дамжуулалт арга нь том хэмжээний загварын мэдлэгийг жижиг, хөнгөн загварт шилжүүлэх үр дүнтэй арга болохыг энэхүү судалгаагаар харуулж байна. Багш загвараас үүсгэсэн өндөр чанарын Q\&A хос өгөгдөл нь сургалтын өгөгдлийг автоматжуулан бий болгох боломжийг олгож, гар аргаар шошголох зардал, хугацааг эрс бууруулах давуу талтай.

\subsection*{Практик ач холбогдол}

LoRA нарийн тохируулгын аргаар сургасан Qwen2.5-7B загвар нь хөнгөн тохируулгатай тул бодит орчинд хэрэгжүүлэхэд илүү тохиромжтой юм. Үүний зэрэгцээ гүйцэтгэлийн алдагдал харьцангуй бага (3--5\%) байгаа нь уг аргын үр ашиг, хэрэгжүүлэх боломжийг баталж байна.

Мөн энэхүү судалгаа нь хяналтын камерын бичлэгээс зодооны зөрчлийг автоматаар илрүүлж, тухайн үйл явдлын талаар Монгол хэл дээр тайлбар үүсгэх боломжийг бүрдүүлж байна. Энэ нь:

\begin{itemize}
    \item Аюулгүй байдлын байгууллагуудын үйл ажиллагааг автоматжуулах.
    \item Хүний оролцоог бууруулж, хариу арга хэмжээний хурдыг нэмэгдүүлэх.
    \item Монгол хэлний тайлбараар хэрэг шалгах үйл явцыг дэмжих.
    \item Шийдвэр гаргалтад нотлох баримт болгон ашиглах.
\end{itemize}

боломжийг бүрдүүлнэ. Цаашид энэхүү аргачлалыг зөвхөн зодооны илрүүлэлтээр хязгаарлахгүйгээр бусад төрлийн хэвийн бус үйлдэл (жишээлбэл, хулгай, осол, үймээн самуун) илрүүлэхэд өргөтгөн ашиглах боломжтой бөгөөд ухаалаг тандалтын систем хөгжүүлэхэд суурь судалгаа болох ач холбогдолтой.

\subsection*{Нийгмийн ач холбогдол}

Монгол улсын нийслэл Улаанбаатар хотод гэмт хэрэг, нийтийн дэг журмын зөрчил жил бүр тохиолддог. Автоматчилсан зодооны илрүүлэлтийн систем нь:

\begin{itemize}
    \item Цагдаагийн хариу арга хэмжээний хугацааг багасгах.
    \item Нотлох баримт бүрдүүлэх үйл явцыг хурдасгах.
    \item Хүний нөөцийн хэрэглээг оновчлох.
    \item Иргэдийн аюулгүй байдлыг дэмжих.
\end{itemize}
